\chapter*{ДҮГНЭЛТ}\addcontentsline{toc}{chapter}{ДҮГНЭЛТ}\markboth{ДҮГНЭЛТ}{}
Энэ ажлаар зогсоолын системд зориулсан Монгол улсын дугаар таних хоёр шатлалт гүн сургалтын системийг бүтээж, нээлттэй эх үүсвэрээс олдсон бүх өгөгдлийг нэгтгэн үнэлгээний протокол тогтоов. Нээлттэй эх үүсвэрийн судалгаагаар GitHub-ийн 68 зураг (тэмдэгт бүрийн хайрцагтай), Roboflow-ийн 984 зураг (хайрцаг, API түлхүүр шаардлагатай), HuggingFace-ийн 8\,823 олон улсын зураг, GitHub репо дотор хадгалагдсан MLUB тэмцээний 7\,394 бодит дугаарын тайралт (текстээр) олдсон бөгөөд монгол дугаарын дизайныг дуурайсан 90\,021 синтетик мөр зураг үүсгэв. Тэмдэгтийн хайрцгаас дугаарын текст сэргээх, дугаарын давхцалгүй хуваалт, сунгасан хоёр мөрт дугаарын хэлбэрийг автоматаар шошголох аргуудыг боловсруулав.

YOLO11n илрүүлэгчийг олон улсын өгөгдөл дээр урьдчилан сургаж монгол зургаар нарийн тохируулахад монгол шалгалтын олонлог дээр mAP$_{50}=0.992$, mAP$_{50\text{--}95}=0.834$ хүрсэн нь зорилтыг хангав. Монгол кирилл 35 үсэг + 10 орны тэмдэгтийн олонлог дээр CRNN+CTC уншигчийг синтетик, GitHub, MLUB гэсэн дөрвөн үе шаттай сургахад 1\,491 бодит дугаар дээр 96.8 хувийн бүрэн зөв уншилт, 0.65 хувийн CER хүрч 95 хувийн зорилтыг давав; хоёр мөрт дугаар 98.4, нэг мөрт 95.8 хувь. Андуурлын шинжилгээгээр алдааны дийлэнх нь Д$\to$Т, Н$\to$Б зэрэг хэлбэрийн ойролцоо хос ба ховор серийн үсэг дээр төвлөрч байгааг тогтоов. Дугаарын дүрэм, 29 аймгийн код, итгэлээр жинлэсэн тэмдэгт тус бүрийн саналд суурилсан олон кадрын баталгаажуулалтыг хэрэгжүүлж, бүхэл урсгал нэг 1080p кадрыг 7.3~мс-д (p95 9.8~мс) боловсруулж 100~мс зорилтыг хангав. Бүтэн дүрстэй 19 дугаарын жижиг тест дээр төгсгөл хүртэлх нарийвчлал 73.7 хувь (9 кадрын саналаар 78.9 хувь) байгаа нь домэйн зөрүү, ховор үсэг, хайрцгийн зөрүүгээс үүдэлтэй бөгөөд зогсоолын камерын бодит өгөгдөл цуглуулах шаардлагыг тодорхой харуулав.

Дараагийн шатанд Roboflow өгөгдлийг нэмж илрүүлэгчийг дахин сургах, зогсоолын камерын бичлэг цуглуулж шошголох, ховор үсгийг тэнцвэржүүлэх, дүрмийн хязгаарлалттай декод, ONNX/TensorRT экспорт, орох-гарах бүртгэл ба хаалга нээх дүрмийг хэрэгжүүлж зогсоолын бүрэн систем болгон хөгжүүлнэ.
